\chapter{Manual de usuario}

Este anexo describe el uso operativo del sistema para el personal del
parqueadero. Está dirigido a quien supervisa el acceso, no a quien administra la
infraestructura: se supone que el sistema ya está desplegado y funcionando.

\section{Requisitos y acceso}

El único requisito del lado del usuario es un navegador actualizado con acceso a
la red local donde se ejecuta la aplicación. En el entorno de desarrollo la
interfaz se sirve en el puerto $5173$ y la interfaz de programación en el $8000$;
el despliegue con contenedores publica la aplicación en el puerto $8080$.

El acceso exige credenciales. El sistema rechaza las solicitudes de consulta que
no presenten una sesión válida, de modo que la pantalla de acceso
(Figura~\ref{fig:if_login}) es el punto de entrada obligatorio.

\begin{enumerate}
    \item Abrir la dirección de la aplicación en el navegador.
    \item Introducir usuario y contraseña provistos por el administrador.
    \item Confirmar. La aplicación abre en el panel de monitoreo.
\end{enumerate}

La sesión se conserva en una cookie que el código del navegador no puede leer, y
se cierra desde el botón de salida del encabezado. Tras varios intentos fallidos,
el servicio aplica una espera antes de admitir un nuevo intento.

\section{Panel de monitoreo}

El panel (Figura~\ref{fig:if_panel}) es la vista que el operador mantiene abierta
durante el turno. La Tabla~\ref{tab:manual_panel} describe cada elemento y qué
acción corresponde cuando el valor no es el esperado.

\begin{longtable}{
        P{0.24}
        P{0.34}
        P{0.32}
    }
    \caption{Elementos del panel de monitoreo}
    \label{tab:manual_panel}\\

    \toprule
    \rowcolor{gray!20}
    \textbf{Elemento} &
    \textbf{Qué indica} &
    \textbf{Qué hacer} \\
    \midrule
    \endfirsthead

    \multicolumn{3}{c}{\textit{Continuación de la Tabla
    \ref{tab:manual_panel}}} \\
    \toprule
    \rowcolor{gray!20}
    \textbf{Elemento} &
    \textbf{Qué indica} &
    \textbf{Qué hacer} \\
    \midrule
    \endhead

    Medidor de ocupación &
    Vehículos dentro del recinto sobre la capacidad total, y el porcentaje
    correspondiente. &
    Contrastar con lo que se observa. Si no coincide, revisar las sesiones
    activas y las sesiones huérfanas. \\
    \midrule

    Ingresos del último día y estadía promedio &
    Cantidad de vehículos confirmados en las últimas $24$ horas y duración media
    de las sesiones cerradas. &
    Sirve de control de tendencia; no requiere acción. \\
    \midrule

    Latencia p95 &
    Tiempo de inferencia por fotograma en el percentil $95$. &
    Un aumento sostenido indica carga en el equipo o un modelo más pesado.
    Consultar la vista de desempeño del modelo. \\
    \midrule

    Cámaras en línea &
    Cámaras que reportan latido, sobre el total configurado. &
    Si el número baja, verificar el proceso de captura del punto de acceso y la
    conectividad de la cámara. \\
    \midrule

    Tarjetas de cámara &
    Puerta asignada, estado y hora del último latido. &
    Es el diagnóstico por cámara: identifica cuál dejó de reportar. \\
    \midrule

    Vehículos activos &
    Matrícula, tipo, color, puerta de ingreso, hora de entrada y tiempo
    transcurrido de cada sesión abierta. &
    Es la lista sobre la que se opera: permite cerrar una sesión manualmente. \\
    \midrule

    Marca de sesión huérfana &
    Sesión cuya permanencia excede el umbral configurado, probablemente sin
    egreso registrado. &
    Buscar el vehículo en el historial y cerrar la sesión con el motivo
    correspondiente. \\
    \midrule

    Panel en vivo &
    Eventos de ingreso, salida y variaciones de ocupación recibidos por el canal
    persistente. &
    No requiere acción; confirma que el sistema está recibiendo eventos. \\
    \midrule

    Ingresos de los últimos siete días &
    Vehículos confirmados por día. &
    Control de tendencia y de estacionalidad semanal. \\
    \bottomrule
    \multicolumn{3}{r}{\footnotesize\textit{Fuente: Elaboración propia}} \\*
\end{longtable}

\section{Búsqueda de incidentes}

Es el procedimiento que se sigue cuando hay que averiguar qué vehículo pasó por el
acceso en un momento determinado, por ejemplo ante un reclamo o un incidente.

\begin{enumerate}
    \item Abrir la vista de incidentes (Figura~\ref{fig:if_incidentes}).
    \item Ajustar el intervalo \emph{Desde} y \emph{Hasta}. Por omisión comprende
        la hora anterior y la siguiente al instante actual; conviene centrarlo en
        el momento recordado.
    \item Elegir un color si se conoce, o dejar \emph{todos}. El filtro es blando
        a propósito: incluye los vehículos cuyo color estimado está próximo al
        pedido, porque el clasificador confunde tonos claros y oscuros y un
        filtro estricto ocultaría al candidato buscado.
    \item Pulsar \emph{Buscar}. Los candidatos aparecen ordenados por cercanía al
        instante consultado, con su recorte, matrícula, tipo, color y confianza.
    \item Identificar el vehículo por la imagen. Los casos sin matrícula
        reconocida permanecen en la lista: son los que requieren revisión visual.
    \item Seleccionar la tarjeta para abrir el detalle.
\end{enumerate}

\section{Detalle del vehículo y cierre manual}

El detalle (Figura~\ref{fig:if_vehiculo}) presenta el recorte de mayor tamaño
disponible, los atributos con su confianza, el identificador interno y la tabla de
sesiones con ingreso, salida, duración, estado y responsable del cierre.

El cierre manual se usa cuando el sistema no pudo cerrar la sesión por sí mismo,
situación habitual cuando el vehículo egresa sin que la matrícula se lea. El
procedimiento es el siguiente:

\begin{enumerate}
    \item Desde el panel, localizar la sesión en la lista de vehículos activos.
    \item Pulsar \emph{Cerrar}.
    \item Escribir el motivo del cierre, o aceptar el texto sugerido. El motivo y
        el usuario quedan registrados junto con el instante de salida.
    \item Confirmar. La ocupación se actualiza y el evento se difunde a los demás
        clientes conectados.
\end{enumerate}

El motivo no es un trámite: es el dato que permite auditar después por qué una
sesión se cerró sin egreso automático. Un cierre sin motivo convierte la
ocupación en un número sin trazabilidad.

\section{Historial de sesiones}

La vista de historial presenta las sesiones registradas, de la más reciente a la
más antigua, con matrícula, tipo, color, ingreso, salida, duración, estado y
responsable del cierre. El campo de búsqueda filtra por matrícula, tipo o color.
La vista solicita las doscientas sesiones más recientes; para intervalos mayores
se consulta el historial paginado del servicio.

\section{Resolución de problemas frecuentes}

\begin{longtable}{
        P{0.22}
        P{0.33}
        P{0.35}
    }
    \caption{Situaciones frecuentes y acción recomendada}
    \label{tab:manual_problemas}\\

    \toprule
    \rowcolor{gray!20}
    \textbf{Situación} &
    \textbf{Causa probable} &
    \textbf{Acción} \\
    \midrule
    \endfirsthead

    \multicolumn{3}{c}{\textit{Continuación de la Tabla
    \ref{tab:manual_problemas}}} \\
    \toprule
    \rowcolor{gray!20}
    \textbf{Situación} &
    \textbf{Causa probable} &
    \textbf{Acción} \\
    \midrule
    \endhead

    Una cámara figura fuera de línea &
    El proceso de captura se detuvo o la cámara no responde. &
    Verificar el proceso de captura del punto de acceso y la conectividad de la
    cámara; el panel toma el último latido recibido. \\
    \midrule

    Vehículo registrado sin matrícula &
    Resolución insuficiente, placa obstruida o lectura por debajo del umbral. &
    Buscar por ventana horaria y color; la imagen del recorte es la evidencia
    disponible. \\
    \midrule

    Sesión abierta sin egreso &
    El vehículo salió sin que se registrara la salida. &
    Verificar en el historial y cerrar manualmente con el motivo. \\
    \midrule

    Color informado que no coincide &
    Confusión entre tonos claros y oscuros, frecuente de noche. &
    Priorizar la imagen sobre el atributo; el color es un auxiliar de búsqueda,
    no una prueba. \\
    \midrule

    Latencia p95 elevada &
    Carga del equipo o ejecución sin motor compilado. &
    Consultar la vista de desempeño del modelo y verificar que el motor
    especializado esté disponible en el equipo. \\
    \midrule

    La consulta responde con error de autorización &
    Sesión vencida o inexistente. &
    Volver a iniciar sesión. \\
    \bottomrule
    \multicolumn{3}{r}{\footnotesize\textit{Fuente: Elaboración propia}} \\*
\end{longtable}

\section{Límites del manual}

Este anexo describe la operación, no la administración del sistema. La creación de
usuarios y la configuración de las cámaras se documentan en la
Sección~\ref{subsec:vistas_admin}; la calibración de zonas, el despliegue de los
modelos y la política de retención de datos corresponden a tareas de
administración que exceden el alcance de este documento.
